\documentclass[11pt]{article}
\usepackage{hyperref}
\usepackage{amsmath,amssymb,amsthm}
\usepackage{graphicx}
\usepackage{booktabs}
\usepackage[margin=1.1in]{geometry}
\usepackage[colorlinks=true,linkcolor=blue,citecolor=blue]{hyperref}

\theoremstyle{definition}
\newtheorem{definition}{Definition}[section]
\newtheorem{example}[definition]{Example}

\theoremstyle{plain}
\newtheorem{theorem}[definition]{Theorem}
\newtheorem{proposition}[definition]{Proposition}
\newtheorem{corollary}[definition]{Corollary}

\theoremstyle{remark}
\newtheorem{remark}[definition]{Remark}
\newtheorem{observation}[definition]{Numerical Observation}
\newtheorem{openproblem}[definition]{Open Problem}

\title{The Bell Angle Transform: A Normalized Framework for\\ Bell-Like Functions and Tail Behavior}
\author{
Zachary A. Coughlin\\
\texttt{\href{mailto:zaaroncou@gmail.com}{zaaroncou@gmail.com}}\\
\texttt{GitHub: \href{https://github.com/zaaroncou-code?tab=repositories}{https://github.com/zaaroncou-code?tab=repositories}}\\
\texttt{LinkedIn: \href{https://www.linkedin.com/in/zachary-c-438814283/?lipi=urn}{https://www.linkedin.com/in/zachary-c-438814283/?lipi=urn}}
}

\date{October 11th, 2026}

\begin{document}

\maketitle

\begin{abstract}
For a symmetric bell-shaped function, normalizing by peak height and half-width at
half-maximum (HWHM) removes the two-parameter ambiguity of location, height, and
width, producing a function $B$ with $B(0)=1$ and $B(1)=\tfrac12$. Applying the
classical alternating (Zak-transform-at-half-frequency) periodization to $B$ and
evaluating at the origin defines a scalar $\varpi=Y\{B\}(0)$, which reduces to the
selected-interval identity $\varpi=2\sum_{k\ge1}[B(2k)-B(2k+1)]$. Under the baseline
hypotheses used here (symmetry, strict decrease, and decay to zero), neither finite
area nor differentiability is required for $\varpi$ to be defined, and $0<\varpi<1$
for every such bell, with every value in this interval attained. Because $\varpi$
depends on $B$ only through its integer values, it is not injective on the full class
of normalized bells and is developed here as a scalar shape statistic useful for
one-sided screening rather than as a distribution identifier; it is likewise not a
universal measure of asymptotic tail order, a distinction made precise by an explicit
example. On matched-half-width mixtures, however, $\varpi$ is exactly affine, giving
closed-form recovery of the mixing parameter for a peak-normalized pseudo-Voigt-type
line-shape family, an application illustrated numerically and, in a small simulation,
recovered from data. A companion representation, the Bell Vector, is also introduced:
applying the classical Faber--Schauder dyadic reconstruction to $B$ after a
compactifying change of variable yields a countable coefficient vector from which $B$
is exactly and uniquely recovered, complementing the scalar $\varpi$ with a complete
description of shape.
\end{abstract}

 
\section{Introduction}
 

Let a bell-like curve be defined as a continuous, symmetric function $f:\mathbb R\to\mathbb R$
(i.e.\ $f(-x)=f(x)$) with $f(x)\to0$ as $x\to \pm \infty$, strictly increasing on
$(-\infty,0]$ and strictly decreasing on $[0,\infty)$, with its maximum at $x=0$. No
differentiability or finite-area condition is imposed in general; these are invoked
later only where specific results need them (a differentiable subclass for the
derivative identity in Section~3.5, and an integrable, sufficiently regular subclass
for the Fourier representation in Section~3.3), while the Range Theorem and the Bell
Critical Constant of Section~5 require only the strict monotonicity and continuity
stated here, used in particular for the HWHM's uniqueness and throughout the range
theorem's proof. Across statistics and other disciplines of math and science, when
analyzing distributions and their respective curves, a special shape that can show up
is the bell-like curve. For instance, the classical bell curve arises from the normal
distribution. It is centered at its mean $\mu$, where the mean, median, and mode
coincide, and its values decrease away from this center according to the squared
distance $\left(x-\mu\right)^2$. Its shape is therefore determined by a central
location and a scale parameter controlling its spread which produces the familiar
smooth and symmetric bell-like curve. Other bell-like curves like t-distribution,
Cauchy and more also show up in other distributions, having special different
features like the size of their tails and more.

With all these special and unique features for bell-like curves, in this paper, I introduce a transform called the Bell Angle Transform
(Definition~\ref{def:bell-angle-transform}) and a scalar quantity $\varpi$ obtained by
evaluating the transform at the origin. Rather than relying on moments, $\varpi$ is built
from differences in the normalized profile across selected intervals. As a result, it
remains defined for bell-like functions (the functions of bell-like curves) whose variance, kurtosis, or even total area
may fail to exist. After normalization by peak height and half-width at half-maximum,
every strictly decreasing bell considered here receives a Bell Angle in the interval
\[
0<\varpi<1,
\]
and every value in this interval is attainable. The quantity therefore provides a common
scale on which bell shapes with very different tail behavior can be compared.

A single number, however, cannot contain the full shape of an arbitrary bell. Indeed,
distinct bell-like functions may have exactly the same value of $\varpi$. To separate
the problem of ranking a bell from that of representing it, I also introduce
the Bell Vector $\vec{w}(B)$. Whereas $\varpi$ compresses a bell into a scalar statistic,
the Bell Vector gives a countable collection of coefficients from which the entire
normalized bell can be reconstructed. For the continuous, symmetric, decaying class
considered in its construction, a fixed universal reconstruction vector
$\vec{\pi}(x)$ satisfies
\[
B(\vec{w},x)=\vec{\pi}(x)\cdot\vec{w}(B)=B(x),
\]
and the Bell Vector is unique relative to this fixed basis. Thus $\varpi$ and
$\vec{w}$ serve different purposes: the former provides a \(1\)-dimensional measure of
selected tail descent, while the latter retains the information necessary to recover the
full normalized profile.

The main focus of this paper remains the scalar Bell Angle and its behavior. I derive
its selected-interval representation, establish its range, calculate reference values
for several common bell families, examine what information it can and cannot determine,
and study its behavior for heavy-tailed functions and matched-width mixtures. The Bell
Vector is developed as an extended and complementary representation rather than as a replacement for
$\varpi$. Together, the normalization, Bell Angle Transform, Bell Angle scalar, and Bell
Vector provide a framework for answering two related questions: how bell-like functions/curves
can be ranked on a common scale, and how their shapes and tail behavior can be represented
and compared without relying solely on classical moment-based quantities.
 
\section{Preliminary Definitions}
 

Before introducing the Bell Angle Transform, several common bell-shaped functions and
mathematical concepts will be referenced throughout this paper.

\subsection{Gaussian Distribution}
The Gaussian, or normal, distribution is given by
\[
f(x)=\frac{1}{\sigma\sqrt{2\pi}}e^{-\frac{(x-\mu)^2}{2\sigma^2}},
\]
where $\mu$ is the mean and $\sigma$ is the standard deviation. Its tails decay rapidly
as $|x|$ increases.

\subsection{Cauchy Distribution}
The standard Cauchy distribution is
\[
f(x)=\frac{1}{\pi(1+x^2)} .
\]
It has much heavier tails than the Gaussian distribution.

\subsection{Student's $t$-Distribution}
The Student's $t$-distribution is a symmetric bell-shaped distribution controlled by its
degrees of freedom $\nu$:
\[
f_\nu(x)=\frac{\Gamma\!\left(\frac{\nu+1}{2}\right)}
{\sqrt{\nu\pi}\,\Gamma\!\left(\frac{\nu}{2}\right)}
\left(1+\frac{x^2}{\nu}\right)^{-\frac{\nu+1}{2}} .
\]
Smaller values of $\nu$ generally correspond to heavier tails.

\subsection{Logistic Function}
The logistic function is $L(x)=\left(1+e^{-x}\right)^{-1}$, and its derivative
\[
L'(x)=\frac{e^{-x}}{(1+e^{-x})^2}
\]
forms a symmetric bell-shaped curve.

\subsection{Super-Gaussian Functions}
A useful family of bell-shaped functions is
\[
f_m(x)=e^{-x^{2m}},\qquad m\geq1 .
\]
As $m$ increases, the center becomes flatter and the sides become steeper.

\subsection{Kurtosis}
Kurtosis is a statistical quantity based on the fourth standardized moment:
\[
\operatorname{Kurt}(X)=\frac{\mathbb{E}\left[(X-\mu)^4\right]}{\sigma^4}.
\]
It is commonly used to describe a distribution's tendency to produce extreme
observations. It is defined only when the fourth moment is finite.

\subsection{Half-Width at Half-Maximum}
For a centered bell-shaped function $f$, the half-width at half-maximum (HWHM) is the
positive value $X$ satisfying $f(X)=f(0)/2$; the full width at half-maximum is
$\operatorname{FWHM}=2X$.

\subsection{Fourier Transform}
Using the convention adopted in this paper,
\[
\widehat{f}(\xi)=\int_{-\infty}^{\infty}f(x)e^{-2\pi i\xi x}\,dx .
\]

\subsection{Periodization}
The periodization of a function is $P_f(x)=\sum_{n\in\mathbb Z}f(x+n)$, and an
\emph{alternating} periodization has the form
\[
P_f^{(-)}(x)=\sum_{n\in\mathbb Z}(-1)^n f(x+n).
\]

\subsection{Zak Transform}
One form of the Zak transform is
\[
Zf(x,\omega)=\sum_{n\in\mathbb Z}f(x+n)e^{-2\pi i n\omega}.
\]
At $\omega=\frac12$ the factor becomes $(-1)^n$, producing an alternating periodization.

\subsection{Lagrange Inversion}
Lagrange inversion is a classical method for obtaining a power-series representation of
the inverse of an analytic function. It is referenced below in connection with the
horizontal normalization of bell-shaped functions.

 
\section{The Bell Angle Framework}
 

This section introduces the main definitions. Some of the machinery behind the
construction (alternating periodization, Fourier representations, Lagrange inversion,
the HWHM width convention) is already established mathematics. What is proposed here
is the way these are combined for \emph{normalized} bell-shaped functions, and the
quantities that result.

\subsection{Universal Bell Normalization}

Let $f$ be a bell-shaped function with maximum at $x=a$. Before applying the transform,
the function is translated and rescaled into a universal form so that differences in
position, height, and width do not affect the comparison between two bells.

Let $X>0$ be the half-width at half-maximum, so that $f(a+X)=f(a)/2$: continuity,
strict decrease on $[a,\infty)$, and decay to $0$ guarantee, by the intermediate value
theorem, that such an $X$ exists and is unique; by symmetry
$f(a-X)=f(a)/2$ as well. Define the normalized bell by
\[
\boxed{\;B(x)=\frac{f(a+Xx)}{f(a)}\;}
\]
This gives
\[
\boxed{\;B(0)=1,\qquad B(-1)=B(1)=\tfrac12 \;}
\]
so every bell entering the transform is placed on the same height and width scale. When
the bell is already centered at $a=0$ this reduces to $B(x)=f(Xx)/f(0)$.

The value $X$ may be obtained numerically from $f(X)=f(0)/2$ (monotonicity guarantees
a unique root, so bisection or Newton iteration is reliable), or, under appropriate
analytic conditions, through Lagrange series inversion applied to
$g(u)=1-f(\sqrt u)/f(0)$, for which the half-height condition reads $g(X^2)=\tfrac12$.
The series route requires $f''(0)\neq0$ and convergence of the inverse series at
$z=\tfrac12$; neither condition is automatic, and flat-top bells such as $e^{-x^4}$ fail
the first.

\begin{remark}[Scale invariance]
$B$ is unchanged if $f$ is replaced by $\alpha f(\beta x)$ for any $\alpha,\beta>0$: the
peak is multiplied by $\alpha$ and the half-width divided by $\beta$, and both cancel.
Removing this two-parameter ambiguity is the entire purpose of the normalization.
\end{remark}

\subsection{The Bell Angle Transform}

\begin{definition}[Bell Angle Transform]\label{def:bell-angle-transform}
Let $f$ be a bell-shaped function and $B$ its normalized form. The \emph{Bell Angle
Transform} is
\[
\boxed{\;Y\{B\}(x)=\sum_{n\in\mathbb{Z}}(-1)^n B(x+n)\;}
\]
Equivalently, in terms of the original centered function,
\[
\boxed{\;Y\{f\}(x)=\frac{1}{f(0)}\sum_{n\in\mathbb{Z}}(-1)^n f\!\left(X(x+n)\right)\;}
\]
\end{definition}

Throughout, the bilateral sum defining $Y\{B\}(x)$ is understood in the paired sense
\[
Y\{B\}(x) = B(x) + \sum_{n=1}^{\infty}(-1)^n\bigl(B(x+n)+B(x-n)\bigr),
\]
i.e.\ as the limit of symmetric partial sums; this is the convention used in every
convergence and linearity argument below, including the mixture argument inside the
proof of Theorem~\ref{thm:range} and Proposition~\ref{prop:affine}.

The transform takes integer-shifted copies of the normalized bell and alternates their
signs:
\[
\cdots+B(x-2)-B(x-1)+B(x)-B(x+1)+B(x+2)-\cdots.
\]

Two points of care. First, the argument in the second form is $X(x+n)$ and not $Xx+n$:
the bell is normalized \emph{first}, and the integer shifts are then applied to the
normalized copies. Second, the series converges for every bell-shaped $f$ by the Leibniz
criterion, since the terms are eventually monotone in $|n|$ and tend to zero. No
integrability assumption is needed, a point that matters in
Section~\ref{sec:zc}.

The alternating sum is the classical Zak transform evaluated at half-frequency. The Bell
Angle Transform therefore does not introduce alternating periodization as new machinery;
it combines that machinery with the normalization above and the quantities below.

\subsection{Bell Angle Fourier Representation}

The alternating structure gives
\[
Y(x+1)=-Y(x),\qquad\text{hence}\qquad Y(x+2)=Y(x).
\]
Suppose, for this subsection only, that $B\in\mathcal S(\mathbb R)$, the Schwartz class
of rapidly decreasing smooth functions; this is a strictly stronger hypothesis than is
needed for $\varpi$ itself, which by Remark~\ref{rem:weak} requires only symmetry,
strict decrease, and decay to zero. Since the Fourier transform is an automorphism of
$\mathcal S(\mathbb R)$, $\widehat B\in\mathcal S(\mathbb R)$ as well, so both $B$ and
$\widehat B$ decay faster than any polynomial; the classical Poisson summation formula,
applied to the modulated function $e^{\pi it}B(t)$, therefore converges absolutely and
uniformly for every $x$ \cite{Grochenig2001}, and yields
\[
\boxed{\;Y\{B\}(x)=2\sum_{k=0}^{\infty}\widehat{B}\!\left(k+\tfrac12\right)
\cos\!\left((2k+1)\pi x\right)\;}
\]
and writing $B_k=2\widehat{B}\!\left(k+\tfrac12\right)$,
\[
\boxed{\;Y\{B\}(x)=\sum_{k=0}^{\infty}B_k\cos\!\left((2k+1)\pi x\right)\;}
\]
Only odd cosine frequencies appear: $\cos(\pi x),\cos(3\pi x),\cos(5\pi x),\dots$, which
is exactly what the antiperiodicity above requires, since even harmonics are $1$-periodic
and cannot occur in a $1$-antiperiodic function. I refer to this as the \emph{Bell Angle
Fourier Representation}. The underlying periodization identity is established
mathematics; it matters here because it leads directly to the scalar defined next.

\subsection{The Bell Angle Scalar}

\begin{definition}[Bell Angle Scalar]\label{def:varpi}
The \emph{Bell Angle} of a normalized bell $B$ is the value of its transform at the
origin,
\[
\boxed{\;\varpi=Y\{B\}(0)\;}
\]
For convenience the symbol $w$ may also be used as shorthand for $\varpi$.
\end{definition}

The word ``angle'' is part of the name of the framework; $\varpi$ is not a geometric
angle measured in degrees or radians.

From the transform,
\[
\varpi=\sum_{n\in\mathbb{Z}}(-1)^n B(n)
=B(0)+2\sum_{n=1}^{\infty}(-1)^n B(n),
\]
using symmetry. The universal normalization $B(0)=1$, $B(1)=\tfrac12$ makes the first two
contributions cancel exactly,
\[
B(0)-2B(1)=1-2\left(\tfrac12\right)=0,
\]
and therefore
\begin{equation}\label{eq:drops}
\boxed{\;\varpi=2\sum_{k=1}^{\infty}\left[B(2k)-B(2k+1)\right]\;}
\end{equation}

This cancellation is the reason the half-maximum convention is used. Had normalization
fixed $B(1)=c$ instead, the leading contribution would be $1-2c$, and $c=\tfrac12$ is the
unique value that removes it. The normalization does not merely fix a scale: it deletes
the central contribution by construction, leaving the behavior of the bell farther out.

\subsection{Bell Intervals and Tail-Drop Strength}

Identity~\eqref{eq:drops} separates the positive side of the bell into alternating unit
intervals. Define the collection of \emph{Bell intervals}
\[
\boxed{\;\mathcal{I}=\{[2,3],[4,5],[6,7],\ldots\}\;}
\]
For a differentiable normalized bell,
$B(2k)-B(2k+1)=\int_{2k}^{2k+1}-B'(x)\,dx$, so
\[
\boxed{\;\varpi=2\sum_{k=1}^{\infty}\int_{2k}^{2k+1}-B'(x)\,dx\;}
\]
This gives the interpretation used throughout: $\varpi$ measures the total amount by
which the normalized bell \emph{descends} across the selected intervals, weighting each
equally. It is a \emph{tail-drop-strength} quantity in this specific, finite-location
sense: a measure of selected-interval descent, not the area under the tails, and
(see the remark following Theorem~\ref{thm:range}) not a universal ordering of
asymptotic tail decay. Two consequences follow, and both are used later: $\varpi$
requires no moments to exist, and $\varpi$ depends on $B$ only through its values at
the integers.

\subsection{Range of the Bell Angle}

\begin{theorem}[Range]\label{thm:range}
For every strictly decreasing normalized bell,
\[
\boxed{\;0<\varpi<1\;}
\]
and both endpoints are approached but never attained. Every value in $(0,1)$ occurs.
\end{theorem}

\begin{proof}
Every term of $\varpi/2=\sum_{k\ge1}[B(2k)-B(2k+1)]$ is positive because
$B(2k)>B(2k+1)$, so $\varpi>0$. For the upper bound, these terms form a subcollection of
a telescoping series over all consecutive integers from $2$ onward:
\[
\frac{\varpi}{2}<\sum_{j=2}^{\infty}\left[B(j)-B(j+1)\right]=B(2)<B(1)=\frac12,
\]
the first inequality strict because the omitted terms $B(3)-B(4),B(5)-B(6),\dots$ are
each positive, and the middle evaluation using $B(j)\to0$.

For sharpness at $0$, the super-Gaussians $B_m(x)=2^{-x^{2m}}$ are normalized bells with
$\varpi\le2B_m(2)=2\cdot2^{-4^m}\to0$ as $m\to\infty$.

For sharpness at $1$, fix a positive integer $k$ and set $\widetilde\lambda_k=k$,
$\widetilde\varepsilon_k=1/k$, $\widetilde b_k=\widetilde\varepsilon_k4^{-k}$,
$\widetilde c_k=(1-\widetilde b_k)/\tanh\widetilde\lambda_k$, and
$\widetilde P_k(x)=\widetilde c_k\tanh(\widetilde\lambda_k x^2)+\widetilde b_kx^{2k}$,
$\widetilde B_k(x)=2^{-\widetilde P_k(x)}$. Since $\tanh\widetilde\lambda_k>0$ and
$\widetilde b_k<1$, $\widetilde c_k>0$; both $\tanh(\widetilde\lambda_kx^2)$ and
$x^{2k}$ are even and strictly increasing in $|x|$, so $\widetilde B_k$ is a
normalized, symmetric, strictly decreasing bell with $\widetilde B_k(0)=1$. Directly,
$\widetilde P_k(1)=\widetilde c_k\tanh\widetilde\lambda_k+\widetilde b_k=(1-\widetilde
b_k)+\widetilde b_k=1$ exactly, for every $k$, so $\widetilde B_k(1)=\tfrac12$ exactly.
At $x=2$: $\widetilde P_k(2)=\widetilde c_k\tanh(4k)+\widetilde\varepsilon_k$; since
$\tanh k,\tanh4k\to1$ exponentially while $\widetilde\varepsilon_k=1/k\to0$ only
polynomially, $\widetilde c_k\tanh4k\to1^-$ and $\widetilde
P_k(2)=1+\tfrac1k+o(\tfrac1k)\to1^+$, hence $\widetilde B_k(2)\to\tfrac12^-$. At $x=3$:
$\widetilde P_k(3)=\widetilde c_k\tanh(9k)+\widetilde\varepsilon_k(9/4)^k\to\infty$,
since $\widetilde\varepsilon_k(9/4)^k=(9/4)^k/k\to\infty$, hence $\widetilde
B_k(3)\to0$. Writing $\widetilde\varpi_k=\varpi(\widetilde B_k)$, the bound $\varpi<1$
already established gives $\widetilde\varpi_k<1$ for every $k$, while
$\widetilde\varpi_k\ge2[\widetilde B_k(2)-\widetilde B_k(3)]\to2(\tfrac12-0)=1$; by
squeezing, $\widetilde\varpi_k\to1$.

Since $0<\varpi<1$ strictly for every such bell, by the inequalities already
established, neither endpoint is ever attained.

To see that every intermediate value is attained, not merely approached, fix
$c\in(0,1)$ and choose $m$ with $\varpi(B_m)<c$ and $k$ with $\widetilde\varpi_k>c$,
using the two limits above. For $\eta\in[0,1]$ set $B_\eta=\eta\widetilde
B_k+(1-\eta)B_m$. Since $\widetilde B_k,B_m$ are even and strictly decreasing on
$(0,\infty)$ with $\widetilde B_k(0)=B_m(0)=1$, $\widetilde B_k(1)=B_m(1)=\tfrac12$, so
is any convex combination: $B_\eta(0)=1$, $B_\eta(1)=\tfrac12$, and $B_\eta$ is
strictly decreasing on $(0,\infty)$ because $\eta\widetilde
B_k(x_1)+(1-\eta)B_m(x_1)>\eta\widetilde B_k(x_2)+(1-\eta)B_m(x_2)$ whenever $0\le
x_1<x_2$. So $B_\eta$ is a normalized bell for every $\eta$. Because the bilateral sum
is taken in the fixed paired convention of Section~3.2 for every bell here, a finite
linear combination of two limits taken along the same partial-sum sequence equals the
linear combination of the limits, giving
\[
\varpi(B_\eta)=\eta\,\widetilde\varpi_k+(1-\eta)\,\varpi(B_m),
\]
an affine, hence continuous, function of $\eta$ (the same fact is recorded
independently, for general matched-width mixtures, as
Propositions~\ref{prop:mixnorm} and~\ref{prop:affine} in Section~6). Since
$B_\eta|_{\eta=0}=B_m$ and $B_\eta|_{\eta=1}=\widetilde B_k$, we have
$\varpi(B_m)<c<\widetilde\varpi_k$, so by the intermediate value theorem applied to the
continuous, affine function $\eta\mapsto\varpi(B_\eta)$ on $[0,1]$, there exists
$\eta^\ast\in(0,1)$ with $\varpi(B_{\eta^\ast})=c$ exactly. As $c\in(0,1)$ was
arbitrary, every value in $(0,1)$ is attained.
\end{proof}

\begin{remark}[Weak hypotheses]\label{rem:weak}
The proof uses only symmetry, strict decrease, decay to zero, and $B(1)=\tfrac12$. No
smoothness and no integrability are required. This is what allows $\varpi$ to be applied
in Section~\ref{sec:zc} to a bell whose total area is infinite.
\end{remark}

\begin{remark}[The second family is not heavy-tailed]
The construction reaching $\varpi\to1$ has \emph{super}-Gaussian decay; it reaches the
ceiling because its descent is postponed, not because it is slow. Consequently no
universal monotone relationship between tail order and $\varpi$ holds, and none is
claimed here.
\end{remark}

\subsection{Reference Values}

\begin{example}[Closed forms]\label{ex:closed}
Three normalized bells have $\varpi$ in closed form:
\[
2^{-x^2}:\;\;\varpi=\vartheta_4\!\left(0,\tfrac12\right)=0.12112420800258050\ldots,
\qquad
\frac{1}{1+x^2}:\;\;\varpi=\frac{\pi}{\sinh\pi}=0.27202905498213316\ldots,
\]
\[
2^{-|x|}:\;\;\varpi=\tfrac13 .
\]
The first is the Jacobi theta value $\vartheta_4(0,q)=\sum_n(-1)^nq^{n^2}$ at
$q=\tfrac12$ \cite{DLMF}; the second follows from the partial-fraction expansion of $\pi\csc(\pi z)$;
the third is a geometric series. All three were confirmed against direct summation.
\end{example}

\begin{observation}[Spectrum of representative bells]\label{obs:spectrum}
\begin{center}
\begin{tabular}{@{}lll@{}}
\toprule
Normalized bell & Description & $\varpi$\\
\midrule
$2^{-x^{6}}$ & super-Gaussian $m=3$ & $1.0842\times10^{-19}$\\
$2^{-x^{4}}$ & super-Gaussian $m=2$ & $3.05176\times10^{-5}$\\
$2^{-x^{2}}$ & Gaussian & $0.121124208003$\\
$\operatorname{sech}(ax)$, $a=\operatorname{arccosh}2$ & $\operatorname{sech}$-type & $0.225052446335$\\
$\left(1+x^{2}\right)^{-1}$ & Cauchy & $0.272029054982$\\
$2^{-|x|}$ & Laplace-type & $1/3$\\
$\left(1+3x^{2}\right)^{-1/2}$ & infinite area & $0.336624306909$\\
\bottomrule
\end{tabular}
\end{center}
The infinite-area row converges slowly ($O(k^{-2})$ term-wise); naive truncation of
$2\sum_{k\le K}[B(2k)-B(2k+1)]$ at $K=5\times10^6$ is already off by $\sim6\times10^{-8}$,
so this entry was recomputed with convergence-accelerated summation.
\end{observation}

 
\section{What the Bell Angle Does Not Determine}\label{sec:notdetermine}
 

Because $\varpi$ is a single number extracted from an infinite-dimensional class, it
cannot identify a bell. This section makes that precise, since it determines how the
statistic may honestly be used.

\begin{proposition}[Level sets admit arbitrarily many parameters]\label{prop:levelsets}
For every $c\in(0,1)$ and every $N\in\mathbb N$, the level set $\{B:\varpi(B)=c\}$
contains a family of normalized bells parametrized injectively by an open subset of
$\mathbb R^N$.
\end{proposition}

\begin{proof}
By \eqref{eq:drops}, $\varpi$ depends on $B$ only through its values at the integers.
Fix a normalized bell $B_0$ with $\varpi(B_0)=c$ and fix $N$. Choose
$2=a_0<a_1<\cdots<a_N<a_{N+1}=3$; since $B_0$ is strictly decreasing,
$B_0(a_0)>\cdots>B_0(a_{N+1})$, so
$\delta:=\tfrac13\min_{0\le j\le N}[B_0(a_j)-B_0(a_{j+1})]>0$. For
$\vec v=(v_1,\ldots,v_N)$ with each $v_j\in(B_0(a_j)-\delta,B_0(a_j)+\delta)$, a direct
check gives $B_0(a_0)>v_1>\cdots>v_N>B_0(a_{N+1})$ with margin at least $\delta$ at
each step. Define $B_{\vec v}$ on $[2,3]$ as the piecewise-linear interpolant through
$(a_0,B_0(a_0)),(a_1,v_1),\ldots,(a_N,v_N),(a_{N+1},B_0(a_{N+1}))$, equal to $B_0$
outside $(2,3)$ and mirrored evenly. This is continuous, strictly decreasing on
$(0,\infty)$, symmetric, and unchanged at every integer, since the deformation is
supported in the open interval $(2,3)$ and its endpoint values are fixed; hence it is a
normalized bell with $\varpi(B_{\vec v})=\varpi(B_0)=c$ for every $\vec v$. Since
$B_{\vec v}(a_j)=v_j$, the map $\vec v\mapsto B_{\vec v}$ is injective, giving the
claimed $N$-parameter family. As $N$ was arbitrary, the level set contains such
families for every $N$.
\end{proof}

\begin{observation}[A collision between two natural shapes]\label{obs:collide}
Non-uniqueness is not confined to engineered deformations. The normalized
$\operatorname{sech}$ bell and the rational bell
$\left[1+(2^{1/q}-1)x^2\right]^{-q}$ at $q^\star=1.80237118914$ agree to roughly 13
significant figures ($\varpi(\mathrm{sech\text{-}type})=0.2250524463350815\ldots$
against $\varpi(B_{q^\star})=0.2250524463351170\ldots$, a residual of order $10^{-14}$
under convergence-accelerated summation), although one tail is exponential and the
other is a power law, differing by more than two orders of magnitude at $x=10$.
\end{observation}

\begin{remark}[Consequences for use]\label{rem:scope}
Two conclusions follow, and they are used in Section~\ref{sec:applications}.
\emph{Identification} of a bell by $\varpi$ requires restricting to a family on which
$\varpi$ is injective. \emph{Rejection}, by contrast, requires no injectivity at all: if
the range of $\varpi$ over a family excludes an observed value, that family cannot fit
the data. One-sided inference remains valid in general.
\end{remark}

 
\section{A Distinguished Value: the Bell Critical Constant}\label{sec:zc}
 

By Remark~\ref{rem:weak}, $\varpi$ is defined for bells whose total area is infinite.
This gives rise to a distinguished value worth recording.

Consider the normalized rational family
\[
B_q(x)=\left[1+\left(2^{1/q}-1\right)x^2\right]^{-q},\qquad q>0,
\]
every member of which satisfies $B_q(0)=1$ and $B_q(1)=\tfrac12$ automatically. Since
$B_q(x)\sim C_q|x|^{-2q}$, the total area is finite exactly when $q>\tfrac12$. At the
boundary the family reaches
\[
B_*(x)=\frac{1}{\sqrt{1+3x^2}},
\]
which still tends to zero but behaves like $(\sqrt3\,|x|)^{-1}$, so its area diverges.

\begin{definition}[Bell Critical Constant]\label{def:zc}
\[
\boxed{\;Z_c=\varpi(B_*)=\sum_{n=-\infty}^{\infty}\frac{(-1)^n}{\sqrt{1+3n^2}}
\approx 0.336624306909156\;}
\]
The subscript $c$ refers to the \emph{critical} behavior of this family.
\end{definition}

The interest of $Z_c$ is that it is a finite shape value attached to a bell for which the
ordinary area statistic has already failed. $Z_c$ is $\varpi$ evaluated at the family's
integrability boundary itself, $B_*=B_{1/2}$; whether $\varpi(B_q)\to Z_c$ continuously
as $q\to\tfrac12^+$ is Open Problem~\ref{op:interchange} and is not assumed here.

\begin{remark}[What $Z_c$ is not]\label{rem:zcnot}
Three clarifications, since each is a natural guess and each is wrong.
\begin{enumerate}
  \item $Z_c$ is not an upper bound for $\varpi$. By Theorem~\ref{thm:range} there
        are normalized bells (integrable ones) with $\varpi$ arbitrarily close to
        $1$, far above $Z_c$.
  \item $Z_c$ is not a gap or discontinuity in the spectrum. By
        Theorem~\ref{thm:range} the set of attainable values is the whole interval
        $(0,1)$, with no gaps anywhere, and $Z_c$ is an ordinary interior point of it.
  \item $Z_c$ is not a universal integrability threshold. Bells with
        $\varpi>Z_c$ may have finite area, and a bell that merely happens to satisfy
        $\varpi=Z_c$ is not thereby critical: criticality is a position within a family,
        not a numerical value.
\end{enumerate}
Other families have their own critical values, and $Z_c$ is the one belonging to this
family. Nothing further is claimed about them here.
\end{remark}

 
\section{Applications}\label{sec:applications}
 

\subsection{A shape statistic that survives where moments do not}

Standard shape indices are moment ratios. Kurtosis requires a finite fourth moment, and
even the variance requires a finite second. For heavy-tailed data these do not merely estimate poorly; they do not exist, and the usual toolkit is
unavailable precisely where shape discrimination is most needed.

Because $\varpi$ is a sum of bounded differences of a bounded function, it is finite for
every bell-shaped function. For the normalized Student's $t$ family with $\nu=2q-1$:

\begin{center}
\begin{tabular}{@{}llll@{}}
\toprule
$\nu$ & $\varpi$ & variance & kurtosis\\
\midrule
$0.5$ & $0.298112$ & infinite & undefined\\
$1$   & $0.272029$ & infinite & undefined\\
$2$   & $0.238608$ & infinite & undefined\\
$4$   & $0.203560$ & finite   & undefined\\
$6$   & $0.185021$ & finite   & finite\\
\bottomrule
\end{tabular}
\end{center}

\subsection{Mixtures: exact recovery of a shape parameter}

Proposition~\ref{prop:levelsets} rules out identification in general. Mixture families
are an exception in which identification is not merely possible but exact, and requires
no conjecture.

\begin{proposition}[Matched-width mixtures are automatically normalized]\label{prop:mixnorm}
Let $B_1,B_2$ be normalized bells and $\eta\in[0,1]$. Then
$B_\eta=\eta B_1+(1-\eta)B_2$ is a normalized bell: it is symmetric, strictly decreasing
on $(0,\infty)$, tends to zero, and satisfies $B_\eta(0)=1$ and $B_\eta(1)=\tfrac12$. In
particular its own HWHM equals $1$, so no renormalization is required.
\end{proposition}

\begin{proof}
Each property is preserved under convex combination, and
$B_\eta(0)=\eta+(1-\eta)=1$, $B_\eta(1)=\eta\cdot\tfrac12+(1-\eta)\cdot\tfrac12=\tfrac12$.
Strict decrease makes the solution of $B_\eta(x)=\tfrac12$ on $(0,\infty)$ unique, hence
equal to $1$.
\end{proof}

\begin{proposition}[$\varpi$ is affine on mixtures]\label{prop:affine}
$\varpi(B_\eta)=\eta\,\varpi(B_1)+(1-\eta)\,\varpi(B_2)$.
\end{proposition}

\begin{proof}
By Proposition~\ref{prop:mixnorm} the mixture is already normalized, so
Definition~\ref{def:varpi} applies to it with no change of scale. Each component series
converges by the Leibniz criterion, and a finite linear combination of convergent series
converges to the corresponding combination of limits.
\end{proof}

\begin{corollary}[Exact inversion]\label{cor:inv}
If $\varpi(B_1)\neq\varpi(B_2)$ then $\eta\mapsto\varpi(B_\eta)$ is affine and strictly
monotone, hence injective on $[0,1]$, with
\[
\boxed{\;\eta=\frac{\varpi-\varpi(B_2)}{\varpi(B_1)-\varpi(B_2)}\;}
\]
Injectivity here is a consequence of affineness, not an assumption.
\end{corollary}

\subsection{Pseudo-Voigt line shapes}

Spectroscopy and powder diffraction routinely model peak shapes by the
\emph{pseudo-Voigt} profile: a convex combination of a Gaussian and a Lorentzian sharing
a common half-width, with $\eta\in[0,1]$ interpolating between instrumental (Gaussian)
and lifetime or strain (Lorentzian) broadening. In normalized coordinates this is
\begin{equation}\label{eq:pv}
B_\eta(x)=\frac{\eta}{1+x^2}+(1-\eta)\,2^{-x^2},
\end{equation}
and by Proposition~\ref{prop:mixnorm} it needs no renormalization at any $\eta$. HWHM is
already the standard width convention in these fields, so the normalization of this paper
is the one practitioners use.

\begin{theorem}[Mixing recovery]\label{thm:pv}
For the profile \eqref{eq:pv},
\[
\varpi(B_\eta)=\vartheta_4\!\left(0,\tfrac12\right)+\eta\,\Delta,
\qquad
\Delta=\frac{\pi}{\sinh\pi}-\vartheta_4\!\left(0,\tfrac12\right)
=0.15090484697955266\ldots,
\]
and therefore
\[
\boxed{\;\eta=\frac{\varpi-\vartheta_4\!\left(0,\tfrac12\right)}
{\dfrac{\pi}{\sinh\pi}-\vartheta_4\!\left(0,\tfrac12\right)}\;}
\]
\end{theorem}

\begin{proof}
Immediate from Example~\ref{ex:closed}, Proposition~\ref{prop:affine} and
Corollary~\ref{cor:inv}.
\end{proof}

\begin{observation}[Verification]\label{obs:pv}
Direct summation against the affine prediction:
\begin{center}
\begin{tabular}{@{}lll@{}}
\toprule
$\eta$ & $\varpi$ computed & affine prediction\\
\midrule
$0$    & $0.12112420800258$ & $0.12112420800258$\\
$0.25$ & $0.15885041974747$ & $0.15885041974747$\\
$0.5$  & $0.19657663149236$ & $0.19657663149236$\\
$0.75$ & $0.23430284323724$ & $0.23430284323724$\\
$1$    & $0.27202905498213$ & $0.27202905498213$\\
\bottomrule
\end{tabular}
\end{center}
Agreement to $14$ digits at every $\eta$; the mixture HWHM was separately confirmed to
equal $1.0$ at $\eta=0.3$ and $\eta=0.7$.
\end{observation}

\begin{remark}[Pseudo-Voigt versus true Voigt]\label{rem:voigt}
The true Voigt profile is a \emph{convolution} of Gaussian and Lorentzian, not a convex
combination, and Proposition~\ref{prop:affine} does not apply to it. The result above
concerns the pseudo-Voigt approximation as defined in \eqref{eq:pv} -- a convex
combination of \emph{peak}-normalized Gaussian and Lorentzian components at matched
half-width. Theorem~\ref{thm:pv} is exact for this peak-normalized, matched-HWHM
parameterization; other treatments in the pseudo-Voigt literature, including the
widely used Thompson--Cox--Hastings line \cite{TCH1987} and its extensions
\cite{IdaAndoToraya2000}, may use different component normalizations or width
conventions, and the mixing coefficient recovered here is not claimed to coincide with
a conventionally defined mixing coefficient without independently verifying and, if
needed, converting between the two. Pseudo-Voigt profiles are widely used in powder
diffraction and related line-profile applications \cite{TCH1987}.
\end{remark}

\subsection{Estimation from data}

The empirical procedure is: form a density estimate $\hat f$; locate its mode (or, if
the symmetry center is known independently, use that) and treat it as the origin,
giving $\hat f(0)$; solve $\hat f(\hat X)=\tfrac12\hat f(0)$ for the half-width; evaluate
$\hat\varpi=\sum_n(-1)^n\hat f(\hat Xn)/\hat f(0)$; and invert by
Theorem~\ref{thm:pv}. The third step requires $\hat f$ only at the points $\hat Xn$, and
the terms decay quickly, so roughly a dozen density evaluations sufficed for the profiles
studied in Observation~\ref{obs:sim} below; no general truncation bound is established.
Both inputs, peak height and half-width, are properties of the bulk of the data rather
than of its extremes, which motivates studying robustness relative to moment estimators;
no such advantage is established here.

\begin{observation}[Exploratory simulation]\label{obs:sim}
Samples of size $n=40{,}000$ were drawn from the density proportional to \eqref{eq:pv},
with $40$ replications per $\eta$. A Gaussian kernel estimate was used, with bandwidth
computed from the central $80\%$ of $|x|$, since Silverman's rule is unusable in the
heavy-tailed regime.
\begin{center}
\begin{tabular}{@{}lll@{}}
\toprule
true $\eta$ & mean $\hat\eta$ & standard deviation\\
\midrule
$0.00$ & $-0.011$ & $0.111$\\
$0.25$ & $0.200$  & $0.119$\\
$0.50$ & $0.451$  & $0.119$\\
$0.75$ & $0.739$  & $0.124$\\
$1.00$ & $0.943$  & $0.132$\\
\bottomrule
\end{tabular}
\end{center}
This is an exploratory illustration rather than a checked, reproducible study: the
random-sampling procedure, exact bandwidth formula, HWHM root-solving method, and
treatment of any $\hat\eta$ estimates outside $[0,1]$ were not preserved in enough
detail for independent reproduction, and the reported bias and standard deviations
should be read accordingly.
\end{observation}
\subsection{Approximate shape screening from empirical curves}
\label{sec:empiricalscreening}

The preceding framework is defined for symmetric normalized bells, whereas a curve
estimated from finite data will generally fail to satisfy these conditions exactly.
Sampling variation, measurement noise, and the density-estimation procedure itself may
introduce small asymmetries even when the underlying distribution is symmetric. The
Bell Angle and Bell Vector may nevertheless be used as approximate shape descriptors
after centering and symmetrizing the estimated curve.

Let $\widehat f$ denote a nonnegative estimated curve obtained from observed data, and
let $\widehat a$ denote its estimated mode or another chosen estimate of its symmetry
center. Define the symmetrized profile
\[
\widehat f_s(u)
=
\frac{
\widehat f(\widehat a+u)
+
\widehat f(\widehat a-u)
}{2},
\qquad
u\ge0.
\]
Extend $\widehat f_s$ evenly to negative $u$. This operation does not assert that the
observed curve is truly symmetric; it extracts the symmetric component about the chosen
center for comparison with the symmetric bell families considered in this paper.

Suppose $\widehat f_s$ is sufficiently bell-like that a positive empirical half-width
$\widehat X$ can be located from
\[
\widehat f_s(\widehat X)
=
\frac12\widehat f_s(0).
\]
The corresponding pseudo-normalized empirical bell is
\[
\boxed{
\widehat B(x)
=
\frac{\widehat f_s(\widehat Xx)}
{\widehat f_s(0)}.
}
\]
By construction,
\[
\widehat B(0)=1,
\qquad
\widehat B(1)=\frac12,
\qquad
\widehat B(-x)=\widehat B(x).
\]
Because $\widehat B$ is estimated from finite data, it need not be perfectly smooth or
strictly decreasing, and the exact theoretical results for normalized bells should not
automatically be interpreted as sampling guarantees for $\widehat B$. Instead,
$\widehat B$ provides an approximate normalized shape to which the descriptors of this
paper may be applied.

A
\subsection{Limitations}

\begin{enumerate}
  \item \textbf{No identification in general.} By Proposition~\ref{prop:levelsets},
        $\varpi$ is a scalar shape statistic useful for one-sided screening, not
        a distribution identifier. Claims of identification must name the family they are
        restricted to.
  \item \textbf{A regime where $\varpi$ is uninformative.} For super-Gaussian profiles
        $2^{-x^{2m}}$ the Bell Angle collapses as $\varpi\le2\cdot2^{-4^m}$: about
        $3\times10^{-5}$ at $m=2$ and $10^{-19}$ at $m=3$. It therefore has no resolving
        power for flat-top order and should not be proposed for beam profiling.
  \item \textbf{No sampling theory.} Consistency, asymptotic variance, bandwidth
        selection, and comparisons against maximum-likelihood or moment-based estimators
        are not established here.
\end{enumerate}

 
\section{The Bell Vector and Universal Reconstruction}\label{sec:bellvector}
 

The construction below is an instance of the classical Faber--Schauder system: the
dyadic hat-function expansion of a continuous function on $[0,1]$, with coefficients
given by the deviation of the true value at each dyadic midpoint from the previous
level's linear interpolant, is due to Faber \cite{Faber1910} and Schauder
\cite{Schauder1927} (see \cite{Semadeni1982} for a modern treatment), and is the
standard first example of an explicit Schauder basis for $C[0,1]$. What is specific
here is the compactifying change of variable $t=|x|/(1+|x|)$ below, which transports
this classical reconstruction from $[0,1]$ to the class of normalized bells on
$\mathbb{R}$, together with the observation, given below, that the HWHM normalization
forces the coarsest coefficient $w_{0,0}$ to vanish identically. The Bell
Reconstruction Theorem is accordingly a bell-specialized restatement of the classical
Faber--Schauder convergence theorem, not an independently new convergence result; the
specialization and the vanishing-coefficient fact are what is specific to the
construction developed here.

The Bell Angle $\varpi$ is a scalar descriptor of a normalized bell, but a single scalar cannot uniquely determine an arbitrary bell. A complete representation may instead be obtained through a countably infinite coefficient vector.
Let
\[
B:\mathbb{R}\to\mathbb{R}
\]
be a continuous, even, normalized bell satisfying
\[
B(0)=1,
\qquad
B(1)=\frac12,
\qquad
\lim_{|x|\to\infty}B(x)=0.
\]
Introduce the compactifying transformation
\[
t=t(x)=\frac{|x|}{1+|x|},
\]
which maps $x\in\mathbb{R}$ into $t\in[0,1)$. Its inverse on $x\ge 0$ is
\[
x=\frac{t}{1-t}.
\]
Define
\[
F(t)=B\!\left(\frac{t}{1-t}\right),
\qquad 0\le t<1,
\]
and extend continuously to $t=1$ by
\[
F(1)=0.
\]
Therefore
\[
F(0)=1,
\qquad
F\!\left(\frac12\right)=\frac12,
\qquad
F(1)=0.
\]
Define the triangular function
\[
\phi(u)=\max\{1-|u|,0\}.
\]
For
\[
n\ge 0,
\qquad
0\le k<2^n,
\]
define the universal reconstruction functions
\[
\pi_{n,k}(x)
=
\phi\!\left(
2^{n+1}\frac{|x|}{1+|x|}
-(2k+1)
\right),
\]
together with
\[
\pi_0(x)=\frac{1}{1+|x|}.
\]
These functions are independent of the particular bell $B$.
For each dyadic interval define
\[
t_L=\frac{k}{2^n},
\qquad
t_M=\frac{2k+1}{2^{n+1}},
\qquad
t_R=\frac{k+1}{2^n}.
\]
The corresponding Bell coefficients are
\[
w_{n,k}
=
F(t_M)
-
\frac{F(t_L)+F(t_R)}{2}.
\]
\begin{definition}[Bell Vector]\label{def:bellvector}
The Bell Vector of $B$ is the countably infinite vector
\[
\boxed{
\vec{w}(B)
=
\left(
1,\,
w_{0,0},\,
w_{1,0},w_{1,1},\,
w_{2,0},w_{2,1},w_{2,2},w_{2,3},
\ldots
\right).
}
\]
\end{definition}
Define the universal reconstruction vector
\[
\boxed{
\vec{\pi}(x)
=
\left(
\pi_0(x),\,
\pi_{0,0}(x),\,
\pi_{1,0}(x),\pi_{1,1}(x),\,
\pi_{2,0}(x),\ldots
\right).
}
\]
The HWHM normalization implies
\[
w_{0,0}
=
F\!\left(\frac12\right)
-
\frac{F(0)+F(1)}{2}
=
\frac12-\frac12
=
0.
\]
\begin{theorem}[Bell Reconstruction Theorem]\label{thm:bellreconstruction}
For every continuous normalized bell satisfying the assumptions above,
\[
\boxed{
B(x)=\vec{\pi}(x)\cdot\vec{w}(B).
}
\]
Equivalently,
\[
\boxed{
B(x)
=
\pi_0(x)
+
\sum_{n=0}^{\infty}
\sum_{k=0}^{2^n-1}
w_{n,k}\pi_{n,k}(x).
}
\]
\end{theorem}
\begin{proof}
For each integer $N\ge 0$, let $F_N$ denote the piecewise-linear interpolation of $F$ on the dyadic grid
\[
\left\{
\frac{k}{2^N}:0\le k\le 2^N
\right\}.
\]
At the coarsest level,
\[
F_0(t)
=
(1-t)F(0)+tF(1).
\]
Since
\[
F(0)=1,
\qquad
F(1)=0,
\]
we obtain
\[
F_0(t)=1-t.
\]
Passing from dyadic level $n$ to $n+1$ introduces the midpoint
\[
t_M=\frac{2k+1}{2^{n+1}}
\]
of each interval
\[
[t_L,t_R]
=
\left[
\frac{k}{2^n},
\frac{k+1}{2^n}
\right].
\]
The previous linear interpolant predicts the midpoint value
\[
\frac{F(t_L)+F(t_R)}{2}.
\]
Hence the required midpoint correction is
\[
w_{n,k}
=
F(t_M)
-
\frac{F(t_L)+F(t_R)}{2}.
\]
Since
\[
\phi\!\left(2^{n+1}t-(2k+1)\right)
\]
equals $1$ at $t=t_M$, vanishes at $t=t_L$ and $t=t_R$, and is linear between these points, the successive interpolants satisfy
\[
F_{n+1}(t)
=
F_n(t)
+
\sum_{k=0}^{2^n-1}
w_{n,k}
\phi\!\left(
2^{n+1}t-(2k+1)
\right).
\]
Iterating yields
\[
F_N(t)
=
1-t
+
\sum_{n=0}^{N-1}
\sum_{k=0}^{2^n-1}
w_{n,k}
\phi\!\left(
2^{n+1}t-(2k+1)
\right).
\]
Because $F$ is continuous on the compact interval $[0,1]$, it is uniformly continuous. The mesh size of the dyadic interpolation tends to zero as $N\to\infty$, so
\[
F_N\to F
\]
uniformly on $[0,1]$.
Therefore
\[
F(t)
=
1-t
+
\sum_{n=0}^{\infty}
\sum_{k=0}^{2^n-1}
w_{n,k}
\phi\!\left(
2^{n+1}t-(2k+1)
\right).
\]
Finally substitute
\[
t=\frac{|x|}{1+|x|}.
\]
Since
\[
1-t=\frac{1}{1+|x|},
\]
we obtain
\[
B(x)
=
\pi_0(x)
+
\sum_{n=0}^{\infty}
\sum_{k=0}^{2^n-1}
w_{n,k}\pi_{n,k}(x).
\]
Hence
\[
\boxed{
B(x)=\vec{\pi}(x)\cdot\vec{w}(B).
}
\]
\end{proof}
\begin{theorem}[Bell Vector Uniqueness]\label{thm:bellvectoruniqueness}
The Bell Vector uniquely determines the bell. That is,
\[
\boxed{
\vec{w}(B_1)=\vec{w}(B_2)
\iff
B_1=B_2.
}
\]
\end{theorem}
\begin{proof}
If
\[
\vec{w}(B_1)=\vec{w}(B_2),
\]
then their reconstruction coefficients are identical. Since the reconstruction vector $\vec{\pi}(x)$ is universal,
\[
\vec{\pi}(x)\cdot\vec{w}(B_1)
=
\vec{\pi}(x)\cdot\vec{w}(B_2)
\]
for every $x$. By the Bell Reconstruction Theorem,
\[
B_1(x)=B_2(x)
\]
for every $x\in\mathbb{R}$.
The reverse implication is immediate from the definition of the coefficients.
\end{proof}
Uniqueness here is relative to the fixed hierarchical basis $\vec\pi(x)$ constructed
above; it is not a claim of uniqueness among all conceivable representations of $B$.

Thus the normalized bell admits the exact correspondence
\[
\boxed{
B(x)
\longleftrightarrow
\vec{w}(B),
}
\]
with universal reconstruction formula
\[
\boxed{
B(\vec{w},x)
=
\vec{\pi}(x)\cdot\vec{w},
}
\]
understood as the convergent series pairing of Theorem~\ref{thm:bellreconstruction},
not an $\ell^2$ inner product; square-summability of $\vec w$ is not established.
The vector $\vec{w}$ therefore contains all bell-specific information, while
$\vec{\pi}(x)$ is universal and independent of the particular bell.

 
\section{Open Questions}
 

\begin{openproblem}[Injectivity on the rational family]\label{op:mono}
Is $q\mapsto\varpi(B_q)$ strictly monotone on $(0,\infty)$? Numerical evidence supports
it. The obvious termwise route fails: writing $D_k(q)=B_q(2k)-B_q(2k+1)$, the terms with
$k\ge2$ decrease in $q$ but $D_1$ does not, rising from about $0.0889$ at $q=0.51$ to a
maximum near $q\approx1.2$ before falling. A proof must handle cancellation between $D_1$
and the remaining terms.
\end{openproblem}

\begin{openproblem}[Limit interchange at the critical bell]\label{op:interchange}
Establish $\lim_{q\to1/2^+}\varpi(B_q)=\varpi(B_*)$ rigorously, justifying the
interchange of the limit with the infinite alternating sum.
\end{openproblem}

\begin{openproblem}[Recovering the bell from the transform]\label{op:coeff}
Does the full sequence $\widehat B\!\left(k+\tfrac12\right)$, $k\ge0$ (equivalently the
whole curve $Y\{B\}(x)$ rather than its value at the origin) determine $B$ within the
class? This is the natural question left open by Proposition~\ref{prop:levelsets}.
\end{openproblem}

\begin{openproblem}[Flat-top normalization]\label{op:flat}
A generalized series-reversion representation of the half-width $X$ when $f''(0)=0$.
\end{openproblem}

\begin{openproblem}[Sampling theory]\label{op:sampling}
Consistency, asymptotic variance, and bandwidth selection for $\hat\varpi$, and
robustness comparisons against moment-based shape estimators.
\end{openproblem}

\begin{openproblem}[True Voigt]\label{op:voigt}
The behavior of $\varpi$ on the convolution family, where affineness fails.
\end{openproblem}

 
\section{Relationship to Known Mathematics}
 

The machinery used here is classical and is not claimed as new. Periodization with
modulation (for $F\in L^1$, the periodization of $e^{2\pi ivt}F(t)$ has Fourier
coefficients $\widehat F(n+v)$) gives the Bell Angle Fourier Representation at
$v=\tfrac12$. The alternating shift sum is the Zak transform at $\omega=\tfrac12$, whose
quasi-periodicity specializes to the antiperiodicity $Y(x+1)=-Y(x)$. Poisson summation
underlies both. Lagrange inversion is classical, as is the partial-fraction evaluation
giving $\pi/\sinh\pi$, and $\vartheta_4$ is a standard Jacobi theta function. Half-width
at half-maximum is a long-established width convention, and the rational family above is
the normalized Student's $t$ / Pearson type VII family under reparameterization
\cite{JohnsonKotzBalakrishnan1995}.

What is proposed as potentially distinctive (as a combination, and not asserted as
novel pending a systematic literature search) is: the universal normalization by peak
height together with half-height width; its pairing with alternating periodization; the
scalar $\varpi$; the selected-interval identity \eqref{eq:drops} that the $B(1)=\tfrac12$
cancellation makes possible; the range theorem with sharp endpoints; the reading of
$\varpi$ as a bell ranking with infinite-dimensional level sets; the constant $Z_c$; and
the exact mixture recovery of Theorem~\ref{thm:pv}.

 
\section{Conclusion}
 

This paper developed two complementary descriptions of a normalized bell-shaped
function. The peak-height/HWHM normalization removes location, height, and width as
free parameters, after which the scalar $\varpi$ compresses the resulting profile into
a single selected-interval descent statistic (called tail-drop strength in this
paper) lying in $(0,1)$, computed from the identity
$\varpi=2\sum_{k\ge1}[B(2k)-B(2k+1)]$ and requiring neither differentiability nor
finite area in general.

This compression is necessarily lossy: because $\varpi$ depends on $B$ only through its
integer values, the level set $\{B:\varpi(B)=c\}$ contains, for every $c$ and every
$N$, families of normalized bells parametrized by arbitrarily many independent
deformation parameters. $\varpi$ should accordingly be read as a scalar shape statistic
for one-sided screening, not a distribution identifier, and, as the collapse of
$\varpi$ for super-Gaussian profiles and the sech/rational-family near-collision both
show, not a universal measure of asymptotic tail order.

Where identification is restricted to a specific family, the picture is different: on
matched-half-width mixtures $\varpi$ is exactly affine, giving closed-form,
family-specific recovery of a mixing parameter, illustrated here for a peak-normalized
pseudo-Voigt-type profile. The Bell Vector then supplies what $\varpi$ alone cannot:
applying the classical Faber--Schauder dyadic reconstruction to the normalized bell
after a compactifying change of variable gives a countable coefficient vector that
recovers $B$ exactly and uniquely, relative to this fixed hierarchical basis.

Several questions raised in the course of this development remain open: whether
$\varpi$ is monotone along the rational family $B_q$; whether
$\lim_{q\to1/2^+}\varpi(B_q)$ equals $\varpi(B_*)$; whether the full transform
$Y\{B\}(x)$, rather than its value at the origin, determines $B$; the flat-top case
where the standard half-width construction fails; sampling theory for the empirical
estimator $\hat\varpi$; and the behavior of $\varpi$ on the true (convolution) Voigt
family, where the affine mixture property does not hold. None of these is resolved
here.

\begin{thebibliography}{99}
\bibitem{Grochenig2001} K. Gr\"ochenig, \emph{Foundations of Time-Frequency Analysis},
  Applied and Numerical Harmonic Analysis, Birkh\"auser Boston, 2001.
\bibitem{Faber1910} G. Faber, Ueber die Orthogonalfunktionen des Herrn Haar,
  \emph{Jahresbericht der Deutschen Mathematiker-Vereinigung} \textbf{19} (1910), 104--112.
\bibitem{Schauder1927} J. Schauder, Zur Theorie stetiger Abbildungen in Funktionalr\"aumen,
  \emph{Mathematische Zeitschrift} \textbf{26} (1927), \url{https://doi.org/10.1007/BF01475440}.
\bibitem{Semadeni1982} Z. Semadeni, \emph{Schauder Bases in Banach Spaces of Continuous
  Functions}, Lecture Notes in Mathematics 918, Springer, 1982.
\bibitem{TCH1987} P. Thompson, D. E. Cox, and J. B. Hastings, Rietveld refinement of
  Debye--Scherrer synchrotron X-ray data from Al$_2$O$_3$,
  \emph{Journal of Applied Crystallography} \textbf{20} (1987), 79--83.
\bibitem{IdaAndoToraya2000} T. Ida, M. Ando, and H. Toraya, Extended pseudo-Voigt
  function for approximating the Voigt profile,
  \emph{Journal of Applied Crystallography} \textbf{33} (2000), 1311--1316.
\bibitem{DLMF} NIST Digital Library of Mathematical Functions, Chapter 20 (Theta
  Functions), \url{https://dlmf.nist.gov/20}.
\bibitem{JohnsonKotzBalakrishnan1995} N. L. Johnson, S. Kotz, and N. Balakrishnan,
  \emph{Continuous Univariate Distributions, Volume 2}, 2nd ed., Wiley, 1995.
\end{thebibliography}

\end{document}
